
\subsubsection{Documentation Framework Design}

Gebru et al. introduced \textit{Datasheets for Datasets}, proposing structured templates documenting dataset motivation, composition, collection processes, preprocessing, and intended uses. The framework drew on electronics industry datasheets to improve transparency and accountability in ML datasets. Mitchell et al. extended this approach to \textit{Model Cards for Model Reporting}, providing analogous documentation for trained models. Pushkarna et al. further refined these ideas with \textit{Data Cards}, emphasizing licensing clarity and preprocessing transparency.

These frameworks address dataset documentation at the design level---providing templates and guidelines for what to document. Boyd et al. validated datasheet effectiveness in controlled settings (N=23 participants), demonstrating improved communication in collaborative ML projects. However, these studies measure framework utility when used, not adoption rates in voluntary practice. This study complements this literature by measuring simulated adoption patterns in sampled repositories.

\subsubsection{Empirical Documentation Studies}

Rondina et al. manually assessed 100 ML datasets and found widespread deficiencies in data collection context (lacking in 25\% of datasets) and preprocessing transparency (lacking in 40\%). Oreamuno et al. analyzed HuggingFace datasets and identified ethics documentation as the weakest component. Gim et al. evaluated FAIR compliance on OpenML, reporting that 0\% of datasets achieve ''Reusable'' status and only 5\% meet ''Findable'' criteria.

These studies establish that documentation gaps exist but share a critical limitation: all use cross-sectional measurements capturing current repository state, not initial release documentation. Without temporal precedence validation, these studies cannot distinguish whether gaps arise from initial non-compliance or documentation degradation over time. This study's temporal measurement at T0+90 days addresses this gap.

\subsubsection{Community Engagement and Software Quality}

Software engineering research has studied how community engagement affects project quality. Mockus and Fielding demonstrated that commit frequency correlates with code quality in large open-source projects. Raymond argued that larger contributor bases improve software outcomes (''many eyeballs make bugs shallow'').

Koch et al. analyzed GitHub ML repositories and found that star counts and contributor diversity correlate with documentation completeness. However, these studies test generic community engagement without isolating specific activity dimensions. This study's mechanism specificity tests---separately measuring commits, contributors, and issue responsiveness---reveal that only sustained commit activity correlates with documentation ($\rho$ = 0.951), while team diversity shows no significant relationship ($\rho$ = 0.028).

\subsubsection{Positioning}

This study differs from prior literature in three ways. First, it provides temporal precedence validation through T0+90 measurement, establishing that documentation gaps exist from initial release. Second, it tests mechanism specificity by isolating commit velocity from contributor count and issue responsiveness. Third, it characterizes component-level heterogeneity, identifying licensing as a critical barrier despite being mechanically simpler than other components.
